\begin{proof}[推论~\ref{cor:quantitative-circuits} 的证明]
我们使用 Williams 的见证电路论证
~\cite[Theorem~6.1]{Wil13}，连同 Jahanjou、Miles 与
Viola 的局部归约~\cite[Theorem~1 and Section~6]{JMV13}。

由 \v{Z}\'ak 的非确定性时间层级~\cite{Zak83}
（取~\cite[Theorem~2.1 in the full version]{Wil13}
的表述形式），固定多带图灵机上的一元语言
$H\in\operatorname{NTIME}(2^n)\setminus\operatorname{NTIME}(2^{n/2})$。
对一切足够大的 $n$，把局部归约以其固定机器作用于 $1^n$，
局部性参数取 $r=2$。时间限额 $T=2^n n^{O(1)}$
涵盖在归约中模拟机器的多对数开销，
给出 $q=n+O(\log n)$ 个子句下标位。
该归约在多项式时间内构造一个解码器，
在每个下标上给出三个变量标号及其符号。
令 $\varphi_n$ 由在全部 $2^q$ 个下标上解码出的子句组成。
它可满足当且仅当 $1^n\in H$。

解码器的 $O(n)$ 个输出位各有一棵常数深度的决策树，
故每个只依赖至多固定数 $w$ 个下标位。
在汉明权重至多 $w$ 的全部串上求值解码器
即可恢复每个支撑。事实上，见证对某一坐标依赖的两个串
可以在支撑之外置零而保持其相异的输出。
同样的求值给出每个支撑上的真值表。
这在多项式时间内构造出 $O(n)$ 个扇入二的门的解码器。

把每个变量标号垫到 $dn$ 位（$d$ 为固定正整数），
并声明全部 $2^{dn}$ 个变量。若对此类 $n$ 有 $|i|=dn$
且 $\varphi_n$ 可满足，令 $W(i)$ 为其按字典序第一个
可满足赋值的第 $i$ 位（变量按其标号排序）。
其他情形一律置 $W(i)=0$。$i$ 的长度确定 $n$。
为判定 $W$，写出 $2^q$ 个子句，问 NP 谕示公式是否可满足；
若可满足，通过询问「当前前缀接零是否仍有可满足扩张」
恢复第一个赋值。至多 $2^{dn}$ 个长度 $2^{O(n)}$
的前缀查询即够，故 $W\in\mathrm E^{\mathrm{NP}}$。

置 $a(t)=\sqrt{\log_2 t}$。检验
定理~\ref{thm:main} 固定密度部分的计数是否为正，
可在时间 $V^b2^{A V/a(V)}$（固定 $A,b>0$）
内判定有 $V$ 个变量、至多 $5V$ 个子句的 $3$-SAT。
对数底的变化被吸收进 $A$。求解器使用多项式空间，
故多带模拟可在每次位访问时扫描其工作区，
成本 $V^{O(1)}$ 被吸收进 $b$。

固定正有理数 $\eta<1/(12A)$，并对 $n\ge2$
置 $s_n=\lfloor\eta n\lfloor a(n)\rfloor\rfloor$。
该阈值可用 $n$ 的二进制长度与
$\lfloor\log_2 n\rfloor$ 的整数平方根计算。
设对所有足够大的 $n$ 有 $C_W(dn)\le s_n$。
为判定 $H$：拒绝非一元输入。在 $1^n$ 上，
猜测一个 $dn$ 个下标位、至多 $s_n$ 个门的电路。
把解码器的三个变量标号馈入该电路的三份副本，
并检验所索引的子句是否被违反。
若检验电路不可满足则接受。
当 $1^n\in H$ 时，在长度 $dn$ 处计算 $W$
的电路给出一个接受的猜测。
当 $1^n\notin H$ 时，每个猜测都违反一个子句。

检验电路有 $q$ 个输入与至多 $3s_n+O(n)$ 个门。
上述编码为每个输入与门各用一个变量，
每门至多四个宽度至多三的子句以强制其真值表，
再用一个单元子句钉住输出。
用未用变量垫齐后得 $U_n=3s_n+Bn$ 个变量与至多
$5U_n$ 个子句（$B$ 为固定整数）。
由 $U_n=(3\eta+o(1))n a(n)$ 与 $a(U_n)\sim a(n)$，
检验耗时 $2^{(3A\eta+o(1))n}$。
猜测与构造电路花费多项式时间，
故对一切足够大的 $n$，$H$
可在时间 $2^{n/2}$ 内非确定性地判定。
有限多个例外长度可以硬编码，
这与 $H$ 的选取矛盾。
于是对无穷多个 $n$ 有 $C_W(dn)>s_n$。
由 $s_n\sim\eta n\sqrt{\log_2 n}$ 与 $d$ 固定，
取 $L=W$ 与任何 $0<c<\eta/(d\sqrt{\log 2})$
即证明 \eqref{eq:quantitative-circuit-lower-bound}。
\end{proof}
